\section{The Prisoner: Free for All}

When I first watched the television series, \emph{The Prisoner}, I assumed it was a warning against the dehumanizing and massification methods of the Western liberal ``democracies''. When I recently saw it again, it dawned on me that it is actually a blueprint for the world to come. Do you remember those ominous, ubiquitous cameras that kept an eye on the village? They are actually a ``good thing'' as the UK has installed tens of thousands of those ominous, ubitiquous cameras.

In the episode ``Free for All', the Prisoner --- \#6 --- is encouraged to run against the village manager \#2 in a staged election. The similarities to the recent Democratic National Convention are hard to miss, right down to the staged and emotion-driven speeches and demonstrations.

In anticipation for the POTUS election of 2008, we note that:

\begin{enumerate}
\item \#6 runs as the agent for ``change''
\item \#2 runs as the ``experienced'' candidate
\item There is only the illusion of choice, since the most fundamental issues are not at stake.
\end{enumerate}


Beyond the specifics, we can note that \#1 is never shown, since the ``true'' ruler of the village is forever unknown and unexposed. All opinions are manufactured, usually willingly by the villagers, but by force and violence, if necessary, as in the case of the Prisoner. There is always the need to prove loyalty, not to any higher standard, but purely to the power holders of the village, whoever they may be.

Although the UK today actually does imprison those who dare to transgress their various and sundry speech laws, the BBC web site find it yourself assures us that the ``Prisoner'' is clearly fiction and utterly unlike the government today. But wouldn't you expect them to say that? What if the Prisoner is the documentary and BBC documentaries are works of fiction? How adolescent!


\flrightit{Posted on 2008-09-01 by Cologero}

